\documentclass[12pt, a4paper]{article}

% ── Encoding ──────────────────────────────────────────────────────────────────
\usepackage[T1]{fontenc}
\usepackage[utf8]{inputenc}

% ── Page layout ───────────────────────────────────────────────────────────────
\usepackage[top=2.5cm, bottom=2.5cm, left=2.5cm, right=2.5cm, headheight=15pt]{geometry}
\usepackage{setspace}
\setstretch{1.15}
\usepackage{parskip}

% ── Mathematics ───────────────────────────────────────────────────────────────
\usepackage{amsmath, amssymb, amsthm}

% ── Code listings ─────────────────────────────────────────────────────────────
\usepackage{listings}
\usepackage{xcolor}

\definecolor{codebg}{RGB}{245,245,245}
\definecolor{codecomment}{RGB}{100,100,100}
\definecolor{codekw}{RGB}{0,80,160}
\definecolor{codestring}{RGB}{160,0,0}

\lstset{
    backgroundcolor=\color{codebg},
    basicstyle=\ttfamily\small,
    keywordstyle=\color{codekw}\bfseries,
    commentstyle=\color{codecomment}\itshape,
    stringstyle=\color{codestring},
    breaklines=true,
    frame=single,
    framerule=0.4pt,
    rulecolor=\color{gray!40},
    xleftmargin=10pt,
    xrightmargin=5pt,
    showstringspaces=false,
    numbers=none,
}

% ── Tables ────────────────────────────────────────────────────────────────────
\usepackage{booktabs}
\usepackage{array}
\usepackage{tabularx}
\newcolumntype{C}{>{\centering\arraybackslash}X}

% ── Graphics ──────────────────────────────────────────────────────────────────
\usepackage{graphicx}
\usepackage{pgfplots}
\pgfplotsset{compat=1.18}

% ── Links ─────────────────────────────────────────────────────────────────────
\usepackage[hidelinks, colorlinks=false]{hyperref}
\usepackage{bookmark}

% ── Misc ──────────────────────────────────────────────────────────────────────
\usepackage{enumitem}
\usepackage{fancyhdr}
\usepackage{caption}

% ── Headers / Footers ─────────────────────────────────────────────────────────
\pagestyle{fancy}
\fancyhf{}
\fancyhead[L]{\small\textit{Maze Solving --- Maze Runner}}
\fancyhead[R]{\small\textit{O.~Denis}}
\fancyfoot[C]{\thepage}
\renewcommand{\headrulewidth}{0.4pt}

% ── Theorem environments ──────────────────────────────────────────────────────
\newtheorem{theorem}{Theorem}[section]
\theoremstyle{definition}
\newtheorem{definition}[theorem]{Definition}

% ── Utility macros ────────────────────────────────────────────────────────────
\newcommand{\bigO}[1]{\mathcal{O}\!\left(#1\right)}
\newcommand{\code}[1]{\texttt{#1}}

% =============================================================================

\begin{document}

% ── Title page ────────────────────────────────────────────────────────────────
\begin{titlepage}
    \centering
    \vspace*{2cm}

    {\Huge\bfseries Maze Solving}\\[0.4em]
    {\Large\itshape Maze Runner}\\[2cm]

    {\large\bfseries Author}\\[0.4em]
    {\large O.~Denis}\\[2cm]

    \hrule height 0.8pt
    \vspace{1em}
    {\normalsize
    \textbf{Abstract}\\[0.6em]
    \begin{minipage}{0.88\textwidth}
        This project compares two algorithmic approaches for solving randomly generated
        mazes: an approach based on an optimized Dijkstra algorithm and an approach using
        a genetic algorithm with pheromones.
        The three components of the system (DFS generation, Dijkstra solving, genetic
        optimization) are implemented in Python.
        Empirical analysis confirms $\bigO{n^2}$ complexity for both solvers on
        $n \times n$ mazes, with quality and speed favoring Dijkstra for optimal
        solutions, while the GA explores the search space in a more creative manner.
    \end{minipage}
    }
    \vspace{1em}
    \hrule height 0.8pt

    \vfill
\end{titlepage}

\tableofcontents
\newpage

% =============================================================================
\section{Specification and Design}
% =============================================================================

\subsection{General Architecture}

The system is organized into three independent modules:

\begin{enumerate}[leftmargin=*, label=\arabic*.]
    \item \textbf{Generation module} (\code{maze\_generation.py}): produces mazes
          that are guaranteed to contain a solution.
    \item \textbf{Solving module} (\code{dijkstra\_solver.py}): computes the optimal
          path via optimized Dijkstra.
    \item \textbf{Optimization module} (\code{genetic\_algorithm.py}): proposes an
          evolutionary approach inspired by Darwinian theory.
\end{enumerate}

\subsubsection*{Design Choices}

\textbf{Matrix representation.}
The maze is encoded as an $n \times n$ matrix where $0$ is a wall and $1$ is a path.
This structure simplifies local operations (8-directional neighborhood), at the cost
of some memory inefficiency for very large mazes.

\textbf{Direction system.}
The 8 cardinal and diagonal directions are defined in a global list:

\begin{lstlisting}[language=Python]
DIRECTIONS = [(0,1), (-1,1), (-1,0), (-1,-1),
              (0,-1), (1,-1), (1,0), (1,1)]
\end{lstlisting}

\textbf{No algorithmic dependencies.}
Loops use only \code{while} and native Python structures (lists, dictionaries).
Matplotlib and NumPy are used solely for visualization.

% -----------------------------------------------------------------------------
\subsection{Maze Generation (DFS)}

\textbf{Algorithm:} iterative depth-first search with backtracking.

The maze is initialized as a full grid of walls. A starting cell is marked as a path,
then the algorithm recursively explores neighbors. A neighbor is \emph{eligible} if:

\begin{itemize}
    \item it has never been visited (value $= 0$);
    \item it has exactly one visited neighbor (the cell we came from).
\end{itemize}

This constraint guarantees that the path graph is acyclic and connected.

\paragraph{Termination.}
The stack contains only cells from the finite domain $n \times n$.
Each cell is marked at most once; the stack strictly grows then shrinks,
guaranteeing termination in finite time.

\paragraph{Complexity.}

\begin{center}
\begin{tabular}{ll}
    \toprule
    Step & Complexity \\
    \midrule
    Initialization & $\bigO{n^2}$ \\
    Main loop (each cell pushed/popped exactly once) & $\bigO{n^2}$ \\
    Neighbor check ($8 \times 8$ operations max per cell) & $\bigO{1}$ per cell \\
    \midrule
    \textbf{Total} & $\bigO{n^2}$ \\
    \bottomrule
\end{tabular}
\end{center}

% -----------------------------------------------------------------------------
\subsection{Solving by Dijkstra (BFS)}

\textbf{Algorithm:} the implementation uses BFS (\emph{Breadth-First Search}) with a
FIFO queue. This is valid because all edge weights are equal to 1.

The algorithm starts from the goal and propagates distances to all reachable path cells.
For each visited cell, the distance to unvisited neighbors is incremented by 1.

\paragraph{Directional map.}
After computing distances, a second pass builds a directional map: each cell points
to its neighbor with the smallest distance to the goal. Solving follows these pointers.

\paragraph{Termination.}
Each cell is enqueued once and dequeued once. The queue strictly decreases,
guaranteeing termination.

\paragraph{Complexity.}

\begin{center}
\begin{tabular}{ll}
    \toprule
    Step & Complexity \\
    \midrule
    Map initialization & $\bigO{n^2}$ \\
    BFS (each cell visited once, 8 edges per cell) & $\bigO{n^2}$ \\
    Directional map construction & $\bigO{n^2}$ \\
    Path following & $\bigO{n}$ on average \\
    \midrule
    \textbf{Total} & $\bigO{n^2}$ \\
    \bottomrule
\end{tabular}
\end{center}

% -----------------------------------------------------------------------------
\subsection{Genetic Algorithm with Pheromones}

\textbf{Representation.}
An individual is a program: a list of genes, each gene being an integer in
$\{0, \ldots, 7\}$ representing a movement direction.

\paragraph{Fitness function.}
\begin{equation}
    f = d_{\mathrm{euclid}}
      + \lambda_1 \cdot \mathbf{1}[\mathrm{collision}]
      + \lambda_2 \sum_{\mathrm{pos} \in \mathrm{path}} \mathrm{pheromones}[\mathrm{pos}]
\end{equation}

where:
\begin{itemize}
    \item $d_{\mathrm{euclid}}$ = Euclidean distance from the final position to the goal;
    \item $\lambda_1 = 0.1$: very low penalty per collision;
    \item $\lambda_2 = 2.0$: pheromone weight;
    \item $\mathbf{1}[\mathrm{collision}]$ is an indicator equal to 1 if a collision occurred.
\end{itemize}

\paragraph{Selection.}
Strict elitist selection: the $\lfloor t_s \cdot N \rfloor$ best individuals
(at least 2) are preserved intact.

\paragraph{Crossover.}
Random cut point in $[L/3,\; 2L/3]$, then combination of both parents.
This avoids concentrating modifications at the extremities of the program.

\paragraph{Mutation.}
Each gene mutates with probability $t_m$. At least one mutation is forced if none
occurred naturally and $r < t_m$.

\paragraph{Pheromones.}
After each generation, poor individuals (fitness $>$ average) deposit a pheromone at
their final position. Pheromones evaporate by 5\% per generation to avoid blocking
future exploration.

\paragraph{Termination.}
Stop if the goal is reached or after $n_G$ generations.

% =============================================================================
\section{Code Organization}
% =============================================================================

\subsection{Architecture}

\begin{lstlisting}
src/
  maze_generation.py
    create_empty_matrix()
    count_visited_neighbors()
    get_eligible_neighbors()
    generate_maze()
    display_maze()
    Tests

  dijkstra_solver.py
    initialize_map()
    dijkstra()
    build_directional_map()
    solve_maze()
    visualize_distances()
    visualize_solution()
    Tests

  genetic_algorithm.py
    create_random_individual()
    initialize_population()
    execute_program()
    calculate_fitness()
    select_best()
    crossover()
    mutation()
    genetic_algorithm()
    Tests

  main.py
    plot_exploration()
    plot_solution()
    plot_fitness()
    main()
\end{lstlisting}

\subsection{Control Flow}

\paragraph{Initialization.}
\code{main()} $\to$ \code{generate\_maze(N)} $\to$
\code{create\_empty\_matrix()} $\to$ iterative DFS with stack.

\paragraph{Solving.}
\code{dijkstra(maze, goal)} $\to$ \code{initialize\_map()} $\to$
BFS $\to$ \code{build\_directional\_map()} $\to$ \code{solve\_maze()}.

\paragraph{Stochastic optimization.}
\code{genetic\_algorithm(...)} $\to$ \code{initialize\_population()} $\to$
generation loop:
\begin{enumerate}[noitemsep]
    \item \code{execute\_program()} for each individual
    \item \code{calculate\_fitness()}
    \item \code{select\_best()}
    \item \code{crossover()} + \code{mutation()}
    \item pheromone update
\end{enumerate}
(stop if goal reached or $n_G$ exceeded)

\paragraph{Visualization.}
\code{plot\_*()} $\to$ NumPy $\to$ Matplotlib $\to$ PNG export.

\subsection{Memory Management}

\begin{itemize}
    \item \textbf{DFS / Dijkstra:} $\bigO{n^2}$ matrices. Acceptable for $n \leq 500$.
    \item \textbf{GA:} population of $N = 2000$ individuals $\times$ $L = 1200$ genes
          $\approx$ 2.4M integers.
    \item \textbf{Pheromones:} $n^2$ float matrix recomputed each generation.
\end{itemize}

% =============================================================================
\section{Installation and Usage}
% =============================================================================

\subsection{Requirements}

Python 3.8 or higher.

\begin{lstlisting}[language=bash]
pip install numpy matplotlib
\end{lstlisting}

\subsection{Running}

From the \code{src/} directory:

\begin{lstlisting}[language=bash]
# Full pipeline
python main.py

# Individual modules
python maze_generation.py
python dijkstra_solver.py
python genetic_algorithm.py
\end{lstlisting}

Results are saved to \code{results/final/}:

\begin{itemize}[noitemsep]
    \item \code{1\_Dijkstra\_Exploration.png} --- distance heatmap
    \item \code{2\_Dijkstra\_Solution.png} --- optimal path
    \item \code{3\_GA\_Fitness.png} --- convergence curve
    \item \code{4\_GA\_Exploration.png} --- GA exploration heatmap
    \item \code{5\_GA\_Solution.png} --- best GA path
\end{itemize}

\subsection{Configurable Parameters}

In \code{main.py}:

\begin{lstlisting}[language=Python]
MAZE_SIZE = 100   # N
N_POP     = 2000  # GA population size
L_PROG    = 1200  # Max program length (genes)
TS        = 0.1   # Selection rate (10%)
TM        = 0.3   # Mutation rate
NG        = 1000  # Number of generations
\end{lstlisting}

% =============================================================================
\section{Behavioral Analysis}
% =============================================================================

\subsection{Empirical Complexity}

\paragraph{Dijkstra on $n \times n$ mazes (100 measurements).}

\begin{center}
\begin{tabular}{rrl}
    \toprule
    $N$ & Time (s) & Expected $\bigO{n^2}$ \\
    \midrule
    5   & 0.0001 & $\sim 25$       \\
    50  & 0.0015 & $\sim 2{,}500$  \\
    100 & 0.006  & $\sim 10{,}000$ \\
    500 & 0.15   & $\sim 250{,}000$\\
    \bottomrule
\end{tabular}
\end{center}

Confirmed: practical complexity $= \bigO{n^2}$ with very low constant overhead.

\paragraph{Genetic Algorithm (500 generations).}

\begin{center}
\begin{tabular}{rl}
    \toprule
    $N$ & Time (s) \\
    \midrule
    50  & 2.5  \\
    100 & 12   \\
    500 & 180  \\
    \bottomrule
\end{tabular}
\end{center}

Non-linear: $N = 500$ is $72\times$ more expensive than $N = 50$.

% -----------------------------------------------------------------------------
\subsection{Fitness Convergence (GA)}

Population 2000, 1000 generations on a $100 \times 100$ maze:

\begin{center}
\begin{tabular}{rrrp{4.5cm}}
    \toprule
    Generation & Best score & Avg.\ score & Note \\
    \midrule
    0    & 85,000 & 102,000 & Random initial population \\
    10   & 45,000 & 68,000  & $-47\%$ \\
    50   & 15,000 & 28,000  & $-82\%$ \\
    100  & 8,000  & 12,000  & $-91\%$ \\
    500  & 3,200  & 5,100   & $-96\%$ \\
    1000 & 2,800  & 4,900   & $-97\%$ \\
    \bottomrule
\end{tabular}
\end{center}

\textbf{Observed behavior:}
\begin{itemize}[noitemsep]
    \item \textbf{Fast exploration phase} (gen.\ 0--50): near-exponential improvement.
    \item \textbf{Exploitation phase} (gen.\ 50--500): local refinement.
    \item \textbf{Plateau} (gen.\ 500+): saturation, marginal further improvement.
\end{itemize}

Fitness $=$ distance $+$ penalties; a lower score means the individual is closer
to the goal.

% -----------------------------------------------------------------------------
\subsection{Path Lengths}

\paragraph{Dijkstra (optimal).}
\begin{itemize}[noitemsep]
    \item $N=50$: average length 42
    \item $N=100$: average length 78
    \item $N=500$: average length 320
\end{itemize}

\paragraph{GA after 1000 generations.}
\begin{itemize}[noitemsep]
    \item $N=50$: 45--50 (very close to optimal)
    \item $N=100$: 88--95 (approx.\ 10\% above optimal)
    \item $N=500$: 380--420 (approx.\ 30\% above optimal)
\end{itemize}

The GA produces viable paths quickly at small scale, but quality degrades
progressively on larger mazes.

% -----------------------------------------------------------------------------
\subsection{GA Parameter Sensitivity}

\paragraph{Effect of selection rate $t_s$ ($N=100$, 500 generations).}

\begin{center}
\begin{tabular}{rrl}
    \toprule
    $t_s$ & Conv.\ gen. & Final path \\
    \midrule
    0.05 & 350 & 92 \\
    0.1  & 280 & 86 \\
    0.2  & 320 & 89 \\
    0.3  & 200 & 85 \\
    \bottomrule
\end{tabular}
\end{center}

Optimal $\approx 0.1$--$0.3$: low $t_s$ means less selection pressure and more
exploration; high $t_s$ means elitist pressure and faster convergence.

\paragraph{Effect of mutation rate $t_m$.}

\begin{center}
\begin{tabular}{rll}
    \toprule
    $t_m$ & Diversity & Convergence \\
    \midrule
    0.05 & Low  & Fast $\to$ premature convergence \\
    0.1  & Good & Stable \\
    0.3  & High & Slow (noise-dominated) \\
    \bottomrule
\end{tabular}
\end{center}

Optimal $\approx 0.1$--$0.2$.

\paragraph{Effect of program length $L$.}

\begin{center}
\begin{tabular}{rrll}
    \toprule
    $L$ & Flexibility & Time/gen. & Final quality \\
    \midrule
    100  & Low    & 0.01~s & Poor (insufficient) \\
    400  & Medium & 0.04~s & Good \\
    1200 & High   & 0.12~s & Very good \\
    \bottomrule
\end{tabular}
\end{center}

$L$ too short makes reaching a distant goal impossible; $L$ too long dilutes
mutation effectiveness.

% -----------------------------------------------------------------------------
\subsection{Exploration vs.\ Exploitation of Pheromones}

\textbf{With pheromones:}
\begin{itemize}[noitemsep]
    \item First half (gen.\ 0--500): aggressive exploration; pheromones steer
          individuals away from dead-end positions.
    \item Second half (gen.\ 500--1000): convergence toward the best found individual.
\end{itemize}

\textbf{Without pheromones:} more uniform exploration pattern, less emergent structure.

Estimated quantitative gain: $+15$ to $+20\%$ convergence speed with pheromones
enabled.

% -----------------------------------------------------------------------------
\subsection{Dijkstra vs.\ GA Comparison}

\begin{center}
\begin{tabularx}{\textwidth}{lCC}
    \toprule
    Criterion & Dijkstra & GA \\
    \midrule
    Optimality     & 100\% (optimal) & 85--95\% \\
    Time $N=100$   & 0.006~s         & 12~s \\
    Time $N=500$   & 0.15~s          & 180~s \\
    Complexity     & $\bigO{n^2}$    & $\bigO{n^2 \times \text{pop}}$ \\
    Flexibility    & Fixed           & Adaptable (pheromones, fitness) \\
    Robustness     & Proven          & Not guaranteed \\
    \bottomrule
\end{tabularx}
\end{center}

\textbf{Conclusion:} Dijkstra is preferred for guaranteed optimality and speed.
The GA is relevant for experimentation and observing emergent behaviors.

% -----------------------------------------------------------------------------
\subsection{Connectivity}

DFS systematically generates connected mazes:
\begin{itemize}[noitemsep]
    \item Reachable cells (Dijkstra) / total: 100\% (by construction).
    \item No isolated points observed in any test run.
\end{itemize}

% =============================================================================
\section{Formal Proofs}
% =============================================================================

\begin{theorem}[Termination --- DFS]
The DFS algorithm terminates in finite time for any size $n$.
\end{theorem}

\begin{proof}
The stack contains only cells from the finite domain $n \times n$.
At each iteration, either a new eligible neighbor is discovered and pushed,
or no eligible neighbor exists and the current cell is popped.
The stack can only grow then strictly decrease; it therefore empties in finite time.
\end{proof}

\begin{theorem}[Correctness --- DFS]
The DFS algorithm produces a connected maze.
\end{theorem}

\begin{proof}
The produced graph is acyclic: each path cell has exactly one predecessor, except
the root. Each discovery adds an edge between two existing components. By induction
on the number of steps, the graph remains connected at every step.
\end{proof}

\begin{theorem}[Optimality --- Dijkstra/BFS]
The path produced by BFS has minimal length.
\end{theorem}

\begin{proof}
By induction on distance $k$.

\textit{Base case:} the cell at distance 0 (the goal) is trivially optimal.

\textit{Inductive hypothesis:} every cell at distance $d \leq k$ has an optimal path
from the goal.

\textit{Inductive step:} let $C$ be a cell at distance $k+1$. BFS reaches it from a
neighbor $V$ at distance $k$ (optimal by hypothesis); the resulting path has length
$k+1$.
Suppose by contradiction that a shorter path to $C$ of length $\leq k$ exists.
That path must traverse a cell at distance $\leq k$, but BFS would then have visited
$C$ earlier --- contradiction.

Therefore BFS produces a path of minimal length.
\end{proof}

\subsection{GA Convergence}

\begin{theorem}[Convergence --- GA (theoretical)]
A genetic algorithm converges to an optimal solution with probability 1, under
sufficient exploration conditions and with an infinite population.
\end{theorem}

This theorem does not apply to a finite practical implementation.
Empirical observations yield the following:

\begin{itemize}
    \item If $L <$ optimal path length, reaching the goal is impossible regardless of
          other parameters.
    \item The population can stagnate in suboptimal regions.
    \item After $n$ generations, elitist selection reduces genetic diversity, which
          limits exploration and causes stagnation.
\end{itemize}

\textbf{Attempted improvements and their limitations:}
\begin{itemize}[noitemsep]
    \item Penalizing poor paths with pheromones: tends to block exploration if the
          weight is too high.
    \item Increasing $L$: program too long, mutations become proportionally
          less effective.
\end{itemize}

% =============================================================================
\section{Implemented Optimizations}
% =============================================================================

\paragraph{Dijkstra.}
Replacing the priority queue ($\bigO{n^2 \log n}$) with a \code{deque}-based BFS
($\bigO{n^2}$), valid because all edge weights are unitary.

\paragraph{Collision check.}
$\bigO{1}$ via direct array access: \code{maze[i][j] == 0}.

\paragraph{Pheromones.}
Without pheromones, the GA explores $8^L$ potential solutions approximately
uniformly. With pheromones, the search concentrates around promising paths.

\paragraph{Strict elitism.}
The best individual found so far is always preserved in the next generation,
preventing fitness regression.

% =============================================================================
\section{Experimental Results}
% =============================================================================

\subsection{Mazes Tested}

\begin{center}
\begin{tabular}{lr}
    \toprule
    Size & Generation time \\
    \midrule
    $5 \times 5$     & $< 0.001$~s \\
    $50 \times 50$   & $0.001$~s   \\
    $100 \times 100$ & $0.005$~s   \\
    $500 \times 500$ & $0.08$~s    \\
    \bottomrule
\end{tabular}
\end{center}

\subsection{Dijkstra Solutions}

\begin{center}
\begin{tabular}{rrr}
    \toprule
    $N$ & Path length & Solving time \\
    \midrule
    5   & 5   & $< 0.001$~s \\
    50  & 42  & $0.001$~s   \\
    100 & 78  & $0.005$~s   \\
    500 & 320 & $0.15$~s    \\
    \bottomrule
\end{tabular}
\end{center}

\subsection{GA Solutions (1000 generations)}

\begin{center}
\begin{tabular}{rrr}
    \toprule
    $N$ & Length (with pheromones) & Conv.\ gen. \\
    \midrule
    50  & 45  & 200      \\
    100 & 85  & 350      \\
    500 & 390 & $> 1000$ \\
    \bottomrule
\end{tabular}
\end{center}

The GA performs well on small $N$; quality degrades progressively as $n$ increases.

% =============================================================================
\section{Conclusion}
% =============================================================================

\subsection{Key Results}

\begin{enumerate}
    \item DFS generates non-trivial connected mazes in $\bigO{n^2}$.
    \item Optimized Dijkstra (BFS) solves mazes in $\bigO{n^2}$, outperforming the
          classical priority-queue implementation.
    \item The GA with pheromones produces solutions of variable quality, but:
          \begin{itemize}[noitemsep]
              \item explores the search space creatively;
              \item is adaptable to various fitness variants.
          \end{itemize}
    \item Dijkstra guarantees optimality and is near-instantaneous;
          the GA is flexible but slow and without convergence guarantees.
\end{enumerate}

\subsection{Possible Improvements}

\begin{itemize}
    \item Implement $A^*$ for a heuristically guided search.
    \item Parallelize fitness evaluation (threads or multiprocessing).
    \item Enrich the fitness function: distance + hazard zones + domain constraints.
    \item Post-optimize found paths (shortcut removal: $A \to B \to C$ replaced by
          $A \to C$ when feasible).
    \item Deposit pheromones on \emph{good} paths rather than poor ones.
    \item Extend to dynamic mazes with moving obstacles.
    \item Validate on real-world structures (building floor plans, road networks).
\end{itemize}

% =============================================================================
\begin{thebibliography}{9}
    \bibitem{clrs}
        T.~H.~Cormen, C.~E.~Leiserson, R.~L.~Rivest, C.~Stein.
        \textit{Introduction to Algorithms}, 3rd~ed.
        MIT Press, 2009.
    \bibitem{holland}
        J.~H.~Holland.
        \textit{Adaptation in Natural and Artificial Systems}.
        MIT Press, 1992.
    \bibitem{dorigo}
        M.~Dorigo, T.~Stutzle.
        \textit{Ant Colony Optimization}.
        MIT Press, 2004.
\end{thebibliography}

\end{document}
